% Generated draft: requires mathematical and source review.
\documentclass[11pt]{article}
\usepackage[margin=25mm]{geometry}
\usepackage{fontspec}
\setmainfont{Latin Modern Roman}
\usepackage{amsmath,amssymb,mathtools,unicode-math}
\setmathfont{Latin Modern Math}
\usepackage{enumitem,needspace,xurl,hyperref}
\setlength{\parindent}{0pt}
\setlength{\parskip}{5pt}
\setlength{\emergencystretch}{3em}
\newcommand{\N}{\mathbb{N}}
\newcommand{\Z}{\mathbb{Z}}
\newcommand{\Q}{\mathbb{Q}}
\newcommand{\R}{\mathbb{R}}
\newcommand{\C}{\mathbb{C}}
\newcommand{\F}{\mathbb{F}}
\newcommand{\A}{\mathbb{A}}
\newcommand{\PP}{\mathbb{P}}
\newcommand{\E}{\mathbb{E}}
\newcommand{\eps}{\varepsilon}
\newcommand{\dd}{\,\mathrm{d}}
\begin{document}
\section*{UnsolvedMath ID: 800844}
\textbf{Source problem code:} catalogue entry
\subsection{Definitions and mathematical statement}
For each positive integer \(n\), let \(\{-1,1\}^{n\times n}\) denote the set of real \(n\times n\) matrices whose entries are all \(-1\) or \(1\). Using the usual sign-discrepancy convention, define
\[
\operatorname{disc}(A):=\min_{x\in\{-1,1\}^{n}}\|Ax\|_{\infty}
=\min_{x\in\{-1,1\}^{n}}\max_{1\leq i\leq n}\left|\sum_{j=1}^{n}a_{ij}x_j\right|.
\]
The problem is to determine the exact value of
\[
D:=\sup_{n\geq 1}\ \sup_{A\in\{-1,1\}^{n\times n}}\frac{\operatorname{disc}(A)}{\sqrt n}.
\]
The supplied record does not explicitly define \(\operatorname{disc}\); the convention displayed above requires confirmation against the source. The restriction \(n\geq1\) makes the normalization well defined.

For context, the supplied excerpt says that an unnamed conjecture implies Spencer's theorem because every column of an \(n\times n\) sign matrix has Euclidean norm \(\sqrt n\). The conjecture and the precise version of Spencer's theorem are not included in the record. The excerpt also attributes to Tim Kunisky, under reference label \texttt{Tim-lowerboundKomlos}, a construction giving the lower bound \(1+\sqrt2\) for the K\'{o}mlos constant \(K\). It describes this as the best current lower bound and says that the constructed matrix is not a sign matrix even after scaling, so the construction does not provide a lower bound in the Spencer setting. These are supplied contextual claims, not independently verified assessments; no precise definition of \(K\) or bibliographic identification of that construction is supplied.
\subsection{Short English statement}
Determine the exact supremum, over all positive sizes and all square sign matrices, of the smallest possible maximum absolute signed row sum divided by the square root of the size. This uses the usual sign-discrepancy convention, which the supplied excerpt leaves undefined. The local registry labels the problem open; this draft does not independently verify that status.
\subsection{Sources}
\begin{itemize}
\item Local UnsolvedMath registry, permanent ID \(800844\), exact source problem code \texttt{catalogue entry}; supplied title: Randomstrasse101 Open Problems (2024) --- Problem 10 (How many deviations are needed afterall?). The registry attributes the problem to Bandeira et al., gives proposed year 2025, and labels its status open.
\item Bandeira et al., \emph{Randomstrasse101: Open Problems of 2024} (2025), Problem 10. Supplied source URL: \url{https://arxiv.org/abs/2504.20539}. The record reports access on 2026-10-16 and extraction by \texttt{arxiv\_latex}; these are supplied provenance details, not a claim of independent source access or verification.
\item Tim Kunisky, cited in the supplied excerpt only by reference label \texttt{Tim-lowerboundKomlos}, for the stated K\'{o}mlos lower-bound construction. No title, publication date, or URL for this reference is supplied, so its bibliographic identity and the accompanying literature assessment remain unverified.
\end{itemize}
\end{document}
